% AUTO-GENERATED by scripts/build_topics.py — do not edit.
% Edit topics/bodies/topic_12.tex or topics/preamble.tex instead.

\begin{konspekt}
\kreq{Примитиви}{създаване на процес, изпълнение на програма, завършване, синхронизация със завършването на процеса-син --- \kref{12.1}.}
\kreq{Права и идентичности}{потребителски идентификатори на процес, групи процеси и сесия --- \kref{12.2}.}
\kreq{Междупроцесна комуникация}{сигнали и програмни канали --- \kref{12.3}.}
\kreq{Междупроцесна комуникация}{IPC пакетът на UNIX System V --- обща памет, семафори, съобщения --- \kref{12.4}.}
\kreq{Схема}{родител--дете --- \kr{12.5.1}, като пълна програма в \kref{12.6}.}
\kreq{Схема}{производител--потребител --- \kref{12.5.2}.}
\kreq{Типове задачи}{задачи, съответни на съдържанието на въпроса.}
\kstatus{схемата родител--дете е дадена като пълна програма; производител--потребител остава схема. Първичните източници липсват в корпуса --- главата подлежи на проверка по литературата.}
\kreq{Литература}{[21], [46].}
\end{konspekt}

\begin{supp}[Статус на източниците]
В наличния корпус липсват първичните материали за този въпрос. Главата следва подробната официална анотация и е сверена с предоставения сравнителен конспект. Преди окончателно академично издание тя трябва да се провери по посочената литература [21] и [46].
\end{supp}

\section{Процес и основни примитиви}\label{s:12.1}

\begin{defn}[Процес]
Процесът е изпълняващ се екземпляр на програма заедно с адресното му пространство, регистрите, отворените файлови дескриптори, идентичността и останалото състояние, което ядрото управлява.
\end{defn}

В Unix създаването и зареждането на програма са отделни операции.

\begin{defn}[Системни примитиви за процеси]
\begin{description}
\item[\texttt{fork}.] Създава процес-син като логическо копие на родителя. И двата процеса продължават след извикването: в сина резултатът е $0$, а в родителя --- PID на сина. При грешка родителят получава $-1$. Реализацията обичайно използва copy-on-write.
\item[\texttt{exec}.] Семейството \texttt{exec} заменя програмния образ на текущия процес с нова програма. При успех не се връща; PID остава същият, а запазването на дескриптори зависи от флага close-on-exec.
\item[\texttt{exit} и \texttt{\_exit}.] Завършват процеса със статус. \texttt{exit} изпълнява регистрираните потребителски обработчици и flush на C потоците, докато \texttt{\_exit} пропуска тази потребителска финализация и предава завършването на ядрото. Ядрото все пак трябва да освободи ресурсите и да затвори дескрипторите, което може да изисква време. След неуспешен \texttt{exec} в дете често се използва \texttt{\_exit}.
\item[\texttt{wait}/\texttt{waitpid}.] Родителят изчаква промяна на състоянието на дете и получава статуса му. Завършило, но неприбрано дете е zombie. Ако родителят не иска да блокира, \texttt{waitpid} позволява подходящи опции.
\end{description}
\end{defn}

Класическият шаблон е: \texttt{fork}; детето подготвя дескрипторите и извиква \texttt{exec}; родителят продължава паралелно или извиква \texttt{waitpid}. Кодът винаги различава трите резултата на \texttt{fork}.

\section{Права, групи и сесии}\label{s:12.2}

\begin{defn}[Идентичност на процес]
Всеки процес има реални и ефективни потребителски и групови идентификатори. Реалните UID/GID описват произхода, а ефективните обичайно участват в проверките за достъп. Допълнителни групи и механизми като set-user-ID променят конкретните правила, затова не е достатъчно да се гледа само PID.
\end{defn}

\begin{defn}[Процесна група]
Процесната група е множество процеси с общ PGID, към което терминалът или друг процес може да адресира сигнал като към едно цяло.
\end{defn}

\begin{defn}[Сесия]
Сесията е множество процесни групи. Обичайно обвивката създава отделна група за всяка задача и използва сесията и управляващия терминал за foreground/background управление.
\end{defn}

Всяка процесна група принадлежи на точно една сесия. \texttt{setpgid} създава или променя принадлежността към процесна група при допустимите ограничения, а \texttt{getpgrp} връща текущата група. \texttt{setsid} създава нова сесия, нова процесна група и прави извикващия процес техен лидер; \texttt{getsid} връща идентификатора на сесията. Тези операции стоят в основата на управлението на задачи от командната обвивка.

\section{Сигнали и програмни канали}\label{s:12.3}

\begin{defn}[Сигнал]
Сигналът е асинхронно известие за събитие. Процесът може да установи обработчик, да остави стандартното действие или за някои сигнали да ги игнорира. Сигналите \texttt{SIGKILL} и \texttt{SIGSTOP} не могат да бъдат прихващани, блокирани или игнорирани. Маска блокира временно доставянето на други сигнали; обработчикът трябва да използва само операции, безопасни в сигнален контекст. \texttt{SIGCHLD} уведомява родителя за промяна в състоянието на дете. При обичайното поведение родителят прибира статуса чрез \texttt{wait}/\texttt{waitpid}; при изрично игнориране на \texttt{SIGCHLD} или флаг \texttt{SA\_NOCLDWAIT} завършилото дете не се оставя като чакащ zombie процес.
\end{defn}

\begin{defn}[Програмен канал]
Програмният канал (pipe) е еднопосочен байтов поток с край за четене и край за запис. \texttt{pipe} връща два файлови дескриптора, които могат да се наследят през \texttt{fork}.
\end{defn}

След \texttt{fork} всеки процес затваря неизползваните краища. Четенето връща край на файл едва когато всички дескриптори за запис са затворени; забравен записващ край може да причини безкрайно блокиране. Записите до определена малка граница са атомарни, но pipe не пази структура на произволно големи съобщения. Именуваният канал FIFO позволява несвързани процеси да отворят общ обект от файловата система.

\section{UNIX System V IPC}\label{s:12.4}

System V IPC идентифицира обекти чрез ключове и системни идентификатори; правата им са отделни от обикновените файлови дескриптори.

\subsection{Обща памет}\label{s:12.4.1}

\begin{defn}[Обща памет]
Общата памет позволява няколко процеса да прикачат един сегмент към адресните си пространства. Типичните операции са създаване/намиране, прикачване, откачване и управление чрез \texttt{shmget}, \texttt{shmat}, \texttt{shmdt} и \texttt{shmctl}. След прикачването обменът е бърз, защото няма копиране през ядрото при всяко обръщение, но общата памет сама по себе си не осигурява взаимно изключване или ред на достъп.
\end{defn}

\subsection{Семафори}\label{s:12.4.2}

\begin{defn}[Семафор]
Семафорът е целочислен синхронизационен обект, променян атомарно. Операцията $P$ (wait/down) изчаква положителна стойност и я намалява, а $V$ (signal/up) я увеличава и може да събуди чакащ процес. Двоичен семафор може да пази критична секция; броящ семафор представя наличен брой еднотипни ресурси. При System V \texttt{semget} създава или намира множество от семафори, \texttt{semop} прилага атомарно вектор от операции, а \texttt{semctl} управлява множеството. Неправилен ред на придобиване може да доведе до взаимно блокиране.
\end{defn}

\subsection{Опашки от съобщения}\label{s:12.4.3}

\begin{defn}[Опашка от съобщения]
Опашката пази отделни съобщения с тип и съдържание. \texttt{msgget} създава или намира опашка, \texttt{msgsnd} изпраща, \texttt{msgrcv} получава, а \texttt{msgctl} управлява жизнения цикъл и параметрите й. Получателят може да избере следващото съобщение или съобщение от определен тип. За разлика от pipe границите между съобщенията се запазват. Капацитетът е ограничен и операциите могат да блокират или да се изпълнят неблокиращо.
\end{defn}

System V обектите могат да останат в ядрото след приключване на процеса, докато не бъдат изрично премахнати. Затова програмата трябва да има ясен собственик и път за почистване.

\section{Синхронизационни схеми}\label{s:12.5}

\subsection{Родител и дете}\label{s:12.5.1}

За изпълнение на програма и получаване на резултата родителят създава дете с \texttt{fork}; детето извиква \texttt{exec}; родителят извиква \texttt{waitpid} и проверява дали детето е завършило нормално и какъв е статусът му. Ако между тях има pipe, краищата се създават преди \texttt{fork}, пренасочват се при нужда с \texttt{dup2}, а всички излишни копия се затварят преди чакането.

\subsection{Производител--потребител}\label{s:12.5.2}

Нека общ кръгов буфер има $N$ места. Използват се броящи семафори
\[
empty=N,\qquad full=0
\]
и двоичен семафор
\[
mutex=1.
\]

Производителят изпълнява $P(empty)$, $P(mutex)$, добавя елемент, после $V(mutex)$ и $V(full)$. Потребителят изпълнява $P(full)$, $P(mutex)$, премахва елемент, после $V(mutex)$ и $V(empty)$. Броящите семафори предотвратяват препълване и четене от празен буфер, а \texttt{mutex} пази индексите и съдържанието. Изчакването за свободно/заето място трябва да е преди заключването на \texttt{mutex}; обратният ред може да задържи критичната секция и да блокира другата страна.

\section{Пълна програма: родител, дете и програмен канал}\label{s:12.6}

Схемата в \S12.5.1 казва \emph{кои} примитиви се използват. Работещата програма
трябва да реши и четирите въпроса, които схемата премълчава: кой край на канала
кой затваря, защо детето излиза с \texttt{\_exit}, какво се прави при
\texttt{EINTR} и как се чете статусът.

\begin{verbatim}
/* Родителят стартира външна програма в дете и чете
   изхода й през канал.
   Компилиране: cc -std=c11 -O2 -o run run.c        */
#include <errno.h>
#include <stdio.h>
#include <sys/wait.h>
#include <unistd.h>

int main(void) {
  int fd[2];                 /* [0] за четене, [1] за писане */
  if (pipe(fd) < 0) {
    perror("pipe");
    return 1;
  }

  pid_t pid = fork();
  if (pid < 0) {
    perror("fork");
    return 1;
  }

  if (pid == 0) {              /* ------- дете ------- */
    /* четящият край не му трябва */
    if (close(fd[0]) < 0) _exit(127);
    if (dup2(fd[1], STDOUT_FILENO) < 0) _exit(127);
    /* дубликатът вече е stdout */
    if (close(fd[1]) < 0) _exit(127);
    execlp("ls", "ls", "-l", (char*)NULL);
    _exit(127);              /* стига се само при неуспех */
  }

  /* ---------- родител ---------- */
  /* без това read никога не вижда край на файл */
  if (close(fd[1]) < 0) {
    perror("close");
    return 1;
  }

  int status = 0;
  char buf[4096];
  for (;;) {
    ssize_t n = read(fd[0], buf, sizeof buf);
    /* детето е затворило канала */
    if (n == 0) break;
    if (n < 0) {
      if (errno == EINTR) continue;
      perror("read");
      status = 1;
      break;
    }
    /* изходът също може да се провали: пълен диск, счупена
       тръба, дескриптор само за четене */
    if (fwrite(buf, 1, (size_t)n, stdout) != (size_t)n) {
      perror("fwrite");
      status = 1;
      break;
    }
  }
  close(fd[0]);

  /* буферът се изпразва тук, докато има кой да съобщи за
     грешката; при изход от main е вече късно */
  if (fflush(stdout) != 0 || ferror(stdout)) {
    perror("stdout");
    status = 1;
  }

  int wstatus;
  while (waitpid(pid, &wstatus, 0) < 0) {
    if (errno != EINTR) {
      perror("waitpid");
      return 1;
    }
  }

  if (WIFSIGNALED(wstatus)) {
    fprintf(stderr, "детето е прекратено от сигнал %d\n",
            WTERMSIG(wstatus));
    return 1;
  }
  if (!WIFEXITED(wstatus)) return 1;

  fprintf(stderr, "детето завърши с код %d\n",
          WEXITSTATUS(wstatus));
  /* успех само ако и детето, и изходът са наред */
  return (WEXITSTATUS(wstatus) == 0 && status == 0) ? 0 : 1;
}
\end{verbatim}

Решенията, които се питат на изпита:

\begin{itemize}
\item \textbf{Затварянето на неизползваните краища не е чистота, а коректност.}
Каналът дава край на файл при четене чак когато \emph{всички} записващи
дескриптори са затворени. Ако родителят задържи \verb|fd[1]|, неговият
\texttt{read} блокира вечно, макар детето отдавна да е завършило.
\item \textbf{Детето излиза с \texttt{\_exit}, не с \texttt{exit}.} След
\texttt{fork} детето има копие на буферите на стандартната библиотека;
\texttt{exit} би ги изчистило втори път и същият текст би се отпечатал два пъти.
\texttt{\_exit} завършва процеса без тези действия.
\item \textbf{\texttt{waitpid} се повтаря при \texttt{EINTR}.} Пристигнал сигнал
прекъсва чакането; това не е грешка и не означава, че детето е завършило.
\item \textbf{Без \texttt{waitpid} детето остава \emph{zombie}}: записът в
таблицата на процесите се пази, докато родителят не прочете статуса му.
\item \texttt{execlp} се връща само при неуспех --- затова редът след него е
изход с код, а не продължение на програмата.
\item \textbf{Записът към изхода се проверява точно както четенето.}
\texttt{fwrite} връща броя записани елементи, а не признак за успех; при пълен
диск, затворена тръба надолу по конвейера или дескриптор, отворен само за
четене, тя записва по-малко и програмата трябва да го забележи. Пропуснато ли
се, програмата връща $0$, след като е загубила част от изхода --- най-лошият
възможен изход, защото извикващият скрипт продължава напред.
\item \textbf{\texttt{fflush} се извиква изрично, преди да се докладва успех.}
Стандартната библиотека буферира; ако изпразването се остави на изхода от
\texttt{main}, грешката настъпва след последната възможност да бъде съобщена и
кодът за връщане е вече определен. \texttt{ferror} улавя и грешка, отбелязана в
по-ранно частично записване.
\item \textbf{Кодът за връщане обединява двете условия.} Успех се докладва само
когато детето е завършило с $0$ \emph{и} целият изход е стигнал до
местоназначението си.
\end{itemize}

\begin{supp}[Обхват]
Схемата производител--потребител (\S12.5.2) е дадена със семафорни операции, но
не като пълна програма: тя изисква и създаване на System V обектите, права,
обработка на \texttt{EINTR} при \texttt{semop} и изрично премахване на обектите
при изход, защото те надживяват процеса. Това остава да се напише.
\end{supp}

\section{Изпитен фокус}\label{s:12.7}

\begin{itemize}
\item Разлика между процес и програма; семантика и резултати на \texttt{fork}, \texttt{exec}, \texttt{exit}/\texttt{\_exit}, \texttt{wait}/\texttt{waitpid}; zombie процес.
\item Реални и ефективни UID/GID, процесни групи, сесии и управляващ терминал.
\item Сигнали и маски; специалният статус на \texttt{SIGKILL} и \texttt{SIGSTOP}; \texttt{SIGCHLD}.
\item Pipe/FIFO, наследяване и правилно затваряне на краищата.
\item Обща памет, семафори и опашки от съобщения в System V IPC, включително основните семейства \texttt{shm*}, \texttt{sem*} и \texttt{msg*}, жизнен цикъл и права.
\item Синхронизация родител--дете и схема производител--потребител с \texttt{empty}, \texttt{full} и \texttt{mutex}.
\item Пълна програма родител--дете през канал (\S12.6): ред на затваряне на краищата, \texttt{dup2}, \texttt{\_exit} вместо \texttt{exit} в детето, повтаряне при \texttt{EINTR} и четене на статуса с \texttt{WIFEXITED}/\texttt{WEXITSTATUS}.
\item Защо незатворен записващ край блокира \texttt{read} завинаги и защо липсващият \texttt{waitpid} оставя zombie процес.
\end{itemize}
